\section{Finite completion by positive interpolation}
\label{sec:certificates}
The remaining degrees require a bound on the whole square, including a
strict comparison arbitrarily close to each maximizing corner.
A finite certificate is a finite collection of data whose verified
inequalities imply that continuous statement. Here its geometric data
are rectangles and their subdivision relations. A terminal rectangle in
the subdivision is called a \emph{leaf}. Its numerical data are outward
enclosures of coefficients or function values: an enclosure $[l,u]$
asserts $l\leq z\leq u$ for the exact quantity $z$, with rounding chosen
to preserve these two inequalities. Its final tests
are strict comparisons with a validated lower bound for the proposed
maximum. The mathematical work in this section explains why passing
these finitely many tests proves an assertion at every point of the
square.

It helps to separate the roles of the objects involved. The target is the
Lebesgue function $\lambda_n$, a sum of absolute values of the cardinal
polynomials. The corner value $\Lambda_c(n)$ is an exact real number,
identified earlier by the cardinal formula. A rational lower endpoint
$L_c\leq\Lambda_c(n)$ is an arithmetic tool for making strict tests.
An upper enclosure on a rectangle describes all points of that rectangle
at once. If that enclosure is strictly below $L_c$, the desired strict
comparison with $\Lambda_c(n)$ follows throughout the rectangle. At a
maximizing corner the proof also retains the exact identity for the
value, so that equality can be classified.

For orientation, consider a continuous function $h$ on a closed
interval. Finitely many observed values leave open what happens between
them. A bound on the interpolation error adds the information that
connects those observations to every intermediate point. For example,
if $|h''|\leq K$ on an interval of width $d$, linear interpolation of
the endpoint values has an error of magnitude at most $Kd^2/8$.
Two endpoint upper bounds and the number $K$ therefore give an upper
bound on the whole interval. The rectangle argument below applies the
same idea in each coordinate. The polynomial argument replaces the
remainder estimate by an exact positive-basis expansion.

The finite degree ranges have distinct convenient representations of
this common principle. For $2\leq n\leq40$, the certificate works directly
with the cardinal polynomials on the physical square. For
$41\leq n\leq319$, it works with the phase envelope and the reflection
defect from Section~\ref{sec:argmax}. The analytic argument already
covers every $n\geq320$. Within the phase range, the analytic far estimate
is available from degree 192, and the analytic corner descent from
degree 256. These thresholds determine which existing comparison is
used on a given region; the final conclusion remains the same strict
maximum and the same two equality points.

Both finite
methods use the same principle: interpolate a smooth signed sum with positive
weights, bound its interpolation error, and only then maximize over signs.
For any finite real vector $(h_\ell)$ the identity
\[
 \sum_\ell|h_\ell|=\max_{\varepsilon_\ell\in\{-1,1\}}
                         \sum_\ell\varepsilon_\ell h_\ell
\]
provides the signed sums. Each is smooth whenever the individual
kernels or cardinals are smooth. The signs realizing the maximum may
change inside a leaf; a bound valid for every complete signed sum
continues to hold across all such changes.

The scalar reason for the displayed identity is that
$\varepsilon_\ell h_\ell\leq|h_\ell|$ for each index, while equality
is achieved by choosing $\varepsilon_\ell$ to be the sign of $h_\ell$.
When $h_\ell=0$, either sign achieves equality. Summing proves both
inequalities needed for the identity. There are finitely many possible
sign vectors, so the maximum is attained. The cardinal values at a
fixed physical point, or the primitive-kernel values at a fixed phase
point, form precisely such a finite vector.

A useful example is $F(x)=|x|$ on $[-1,1]$. Its signed family consists
of the two smooth functions $x$ and $-x$, and $F$ is their maximum.
At the origin the two functions meet and their maximizing roles change.
The graph has a cusp there, although each member of the family is a
polynomial. More complicated sums of absolute values can have many
such changes. A single estimate valid for every member survives the
maximum, including all crossing points. This is the reason the
interpolation lemma is stated for a finite maximum of smooth functions.

The full sign vector is fixed while differentiating and interpolating.
At the end it may be chosen separately at every evaluation point.
For instance, if $|f_\sigma|\leq M$ and a frequency estimate gives
$|(f_\sigma)_{aa}|\leq K$ for every fixed $\sigma$, the same $K$
can be used before taking the maximum. Its size is governed by the
whole signed sum. This will keep the derivative allowance independent
of the number of cardinal functions or board points. A separate
estimate for each component would often lose this useful feature.

There are two kinds of positivity in the next lemma. Bilinear
interpolation takes a weighted average of four vertex values, with
nonnegative weights. The tensor Bernstein expansion takes a weighted
average of coefficient values, again with nonnegative weights. Both
families of weights sum to one. Consequently upper bounds can replace
the interpolated values or coefficients without reversing the
inequality. This order-preserving property is what makes outward
arithmetic compatible with the continuous proof.

The Bernstein basis mentioned in the lemma consists of the polynomials
$B_i^n(s)=\binom ni s^i(1-s)^{n-i}$ for $0\leq i\leq n$ and
$0\leq s\leq1$. A tensor basis in two variables consists of their
products $B_i^n(s)B_j^n(t)$. The normalized coordinates $s,t$ are the
affine coordinates carrying the rectangle to $[0,1]^2$. The phrase
``tensor Bernstein coefficients'' means the coefficients of a polynomial
in this product basis. The parameters $\omega_a,\omega_b$ in the
smooth interpolation estimate record the scale of the second derivative
in each coordinate; for the phase application they will be $\pi$ and
$\pi k$.

\begin{lemma}[Positive interpolation bounds]\label{cert:positive}
Let $F(z)=\max_{\sigma}f_\sigma(z)$, where the maximum is finite.
On a rectangle $R=[a_0,a_1]\times[b_0,b_1]$, suppose
$|(f_\sigma)_{aa}|\leq M\omega_a^2$ and
$|(f_\sigma)_{bb}|\leq M\omega_b^2$ for every $\sigma$.
If $W_z\geq F(z)$ at its vertices, then throughout $R$,
\begin{equation}\label{cert:peano}
 F(a,b)\leq I_RW(a,b)+\frac M2\{
 \omega_a^2(a-a_0)(a_1-a)+\omega_b^2(b-b_0)(b_1-b)\}.
\end{equation}
Here $I_R$ is bilinear interpolation. For polynomial cardinals with tensor
Bernstein coefficients $c_{\ell,ij}$ on $R$, the analogous bound is
\begin{equation}\label{cert:bernstein}
 \sum_\ell|L_\ell(z)|\leq
 \sum_{i,j}\left(\sum_\ell|c_{\ell,ij}|\right)B_i^n(s)B_j^n(t),
 \qquad (s,t)\in[0,1]^2.
\end{equation}
\end{lemma}
\begin{proof}
The one-dimensional linear interpolation remainder is an integral against
a Peano kernel of one sign, whose mass is $(a-a_0)(a_1-a)/2$.
Let $I_a,I_b$ denote the two linear interpolation operators. The identity
\[
 f_\sigma-I_aI_bf_\sigma
   =(f_\sigma-I_af_\sigma)+I_a(f_\sigma-I_bf_\sigma)
\]
separates the two errors. The first is bounded by
$M\omega_a^2(a-a_0)(a_1-a)/2$; positivity of $I_a$ bounds the second
by $M\omega_b^2(b-b_0)(b_1-b)/2$. At the four vertices,
$f_\sigma\leq F\leq W$, so positivity also gives
$I_aI_bf_\sigma\leq I_RW$. This proves \eqref{cert:peano} for
every signed sum, with an error independent of $\sigma$, and taking
their maximum proves the assertion.
Equation~\eqref{cert:bernstein} follows from the same positivity and
partition of unity of the Bernstein basis, after the triangle inequality.
\end{proof}

The one-variable remainder used in the proof can be derived directly
from the fundamental theorem of calculus. Let $u<v$, $x\in[u,v]$,
and let $h$ have a continuous second derivative. Write
\[
 (Ih)(x)=\frac{v-x}{v-u}h(u)+\frac{x-u}{v-u}h(v).
\]
Twice integrating $h''$ from $u$ gives
\[
 \begin{aligned}
 h(x)&=h(u)+(x-u)h'(u)+\int_u^x(x-t)h''(t)\,dt,\\
 h(v)&=h(u)+(v-u)h'(u)+\int_u^v(v-t)h''(t)\,dt.
 \end{aligned}
\]
Substitute the second identity into $Ih$ and subtract. The constant and
linear parts cancel. Splitting the integral at $x$ leaves
\[
 h(x)-(Ih)(x)=-\int_u^v K_x(t)h''(t)\,dt,
 \qquad
 K_x(t)=
 \begin{cases}
 \dfrac{(v-x)(t-u)}{v-u},&u\leq t\leq x,\\[4pt]
 \dfrac{(x-u)(v-t)}{v-u},&x\leq t\leq v.
 \end{cases}
\]
Every factor defining $K_x$ is nonnegative. Its graph consists of two
line segments meeting at $x$, and the areas of the resulting two
triangles add to
\[
 \int_u^vK_x(t)\,dt=\frac{(x-u)(v-x)}2.
\]
Thus $|h''|\leq K$ implies
$h(x)\leq(Ih)(x)+K(x-u)(v-x)/2$. A convex function has
$h''\geq0$, in which case the signed remainder itself is nonpositive,
as expected from the secant-line picture. The upper error estimate
uses the absolute bound and remains valid for either curvature sign.
At an endpoint the kernel vanishes and interpolation is exact.

For a concrete check take $h(t)=t^2$ on $[0,1]$. The interpolant is
$x$, and the error is $x^2-x=-x(1-x)$. The second derivative is two,
so the integral formula gives exactly that error. For $h(t)=-t^2$
the sign changes and the upper error allowance is attained. This
example verifies the sign and the factor $1/2$ in the kernel mass.
Those same factors enter the two-dimensional certificate.

To see the weights in two dimensions, put
$s=(a-a_0)/(a_1-a_0)$ and $t=(b-b_0)/(b_1-b_0)$ for a rectangle
of positive widths. The bilinear interpolant of vertex data is
\[
 \begin{aligned}
 I_RW={}&(1-s)(1-t)W_{00}+s(1-t)W_{10}\\
        &+(1-s)tW_{01}+stW_{11}.
 \end{aligned}
\]
The four weights are nonnegative, and their sum is
$((1-s)+s)((1-t)+t)=1$. Hence $I_RW$ lies between the smallest
and largest vertex datum. Increasing any datum increases or preserves
the interpolant. The operators $I_a$ and $I_b$ act on separate
coordinates, and either order gives this same formula.

The error decomposition in the proof has a practical advantage. Its
first term uses only a second derivative in $a$; its second uses only
a second derivative in $b$. The positivity of $I_a$ transfers the
second error estimate from the two $a$ endpoints to every interior
$a$. In particular the second-coordinate factor
$(b-b_0)(b_1-b)$ is constant as $I_a$ acts. This produces exactly
the sum in \eqref{cert:peano}. No mixed derivative estimate is
required for this argument. Each bound is applied to the smooth
function $f_\sigma$; the resulting common allowance can then pass
through the finite maximum defining $F$.

Upper vertex data can have different amounts of slack. Suppose the
exact vertex values of $F$ are $F_{ij}$ and the certified upper data
are $W_{ij}=F_{ij}+e_{ij}$, with $e_{ij}\geq0$. Their extra contribution
to the bilinear interpolant is the weighted average
\[
 (1-s)(1-t)e_{00}+s(1-t)e_{10}+(1-s)te_{01}+ste_{11}.
\]
It is nonnegative and at most the largest $e_{ij}$. Thus outward
rounding has a visible and controlled effect on the bound. Subdivision
can reduce the curvature allowance because its coordinate products
scale quadratically with the widths. Refining the numerical enclosures
can reduce the data slack. The successful leaf inequalities certify
that the final available slack suffices for the desired strict test.

For the polynomial half of the lemma, define the degree-$n$ Bernstein
basis on $[0,1]$ by
\[
 B_i^n(s)=\binom ni s^i(1-s)^{n-i},\qquad 0\leq i\leq n.
\]
The endpoint conventions follow from polynomial evaluation, so
$B_0^n(0)=B_n^n(1)=1$ and all other weights at the corresponding
endpoint are zero. The binomial theorem proves
$\sum_iB_i^n(s)=(s+(1-s))^n=1$. Taking products gives the same
properties on the unit square. If
$L_\ell=\sum_{i,j}c_{\ell,ij}B_i^n(s)B_j^n(t)$, the triangle
inequality gives
\[
 |L_\ell|\leq\sum_{i,j}|c_{\ell,ij}|B_i^n(s)B_j^n(t).
\]
Sum over $\ell$ and interchange finite sums to obtain
\eqref{cert:bernstein}. Each coefficient absolute sum is an upper
control value, and the basis supplies the convex weights. The
coefficients are coordinates in a polynomial basis. Except at a
vertex, they need not equal values of the polynomial at grid points.

For example, $p(s)=s(1-s)$ has degree-two Bernstein coefficients
$0,1/2,0$, since $B_1^2(s)=2s(1-s)$. Positivity bounds $p$ by
$1/2$ on the interval; its actual maximum is $1/4$. The gap simply
measures how conservative this coefficient bound can be. On smaller
intervals the converted coefficients give new local bounds. Exact
subdivision computes those local coefficients while preserving the
polynomial identity. The certificate uses whichever local bounds
actually pass its strict comparisons.

\subsection{One polynomial certificate for the small degrees}
For $2\leq n\leq40$ the direct certificate uses the corrected cardinal
kernel of Lemma~\ref{lem:cardinal-kernel}. Outward rational enclosures of
its coefficients are converted to tensor Bernstein form on the square.
Binary midpoint subdivision uses exact de Casteljau averages. On each leaf
all coefficient absolute sums in \eqref{cert:bernstein} must be strictly
less than a validated lower endpoint $L_c\leq\Lambda_c(n)$.
A leaf containing a designated corner exempts only its vertex coefficient:
that coefficient is exactly the polynomial value there, with absolute sum
$\Lambda_c(n)$; every other coefficient sum must be less than $L_c$.

The affine coordinates on a leaf are the $s,t$ used above, each ranging
from zero to one. Although the Padua interpolation space has a total
degree restriction, every cardinal also has degree at most $n$ in each
coordinate. It therefore has a representation in the tensor degree
$(n,n)$ basis. Some high-degree information is redundant for this
purpose, which is harmless: the identity being bounded is still the
exact corrected cardinal polynomial. The tensor basis is convenient
because positivity and subdivision both act one coordinate at a time.

The lower endpoint $L_c$ deserves careful attention. To prove that an
ordinary leaf lies strictly below the exact corner value, an upper
coefficient sum $U_{ij}$ is checked by the rational inequality
$U_{ij}<L_c$. Since the exact coefficient sum is at most $U_{ij}$ and
$L_c\leq\Lambda_c(n)$, transitivity gives the desired comparison.
There is a positive rational margin $L_c-U_{ij}$. A checker can verify
its sign using integer arithmetic on the numerator after putting the
two rational numbers over a common positive denominator.

A designated corner leaf uses one additional exact fact: evaluating a
polynomial at a vertex selects precisely its vertex Bernstein
coefficient. For the lower-left vertex, for example, $s=t=0$, so all
tensor basis terms vanish except $B_0^n(0)B_0^n(0)=1$. Thus the true
coefficient $c_{\ell,00}$ equals $L_\ell$ at that corner. Its absolute
sum equals the exact number $C=\Lambda_c(n)$. This equality is retained
symbolically. Outward intervals used for the remaining coefficients
supply their upper bounds. The equality at the exceptional vertex and
the strict bounds at all the others have different roles in the same
positive weighted sum.
Write $C=\Lambda_c(n)$, let $w_*$ be the exceptional
vertex weight, and let $D_R$ be the largest validated coefficient sum
among the other indices. The corner predicate gives $D_R<L_c\leq C$,
so the upper bound on this leaf is
\[
 Cw_*+D_R(1-w_*)=C-(C-D_R)(1-w_*).
\]
Here $w_*$ is a product of endpoint Bernstein weights, and $w_*<1$
at every point other than the designated vertex. This proves the
required strict comparison however close the point is to that vertex.
At the vertex itself the cardinal formula supplies the exact value $C$.

For clarity, if the designated vertex is the lower-left one then
$w_*=(1-s)^n(1-t)^n$. For $n\geq1$, both factors lie in $[0,1]$.
Their product is one exactly when $s=t=0$. For any other point at
least one factor is strictly smaller than one, so $1-w_*>0$. At
another vertex, the factors are replaced by the appropriate choices
$s^n,(1-s)^n,t^n,(1-t)^n$, with the identical conclusion. This
argument includes every point on an edge of the rectangle.

A toy coefficient pattern illustrates the strictness. Suppose a
positive-basis upper sum has exceptional value $C=2$ at the lower-left
vertex and every other coefficient sum is at most $D_R=3/2$. The bound
is $2-(1-w_*)/2$. At each point other than the lower-left vertex it
is strictly below two, even when that point is extremely close to the
vertex. The margin tends continuously to zero as the point approaches
the vertex. The actual certificate supplies its own degree-dependent
values $C,D_R$; the logical use of the margin is exactly the same.

This explains how a finite inequality can establish strictness on a
punctured neighborhood. One finite positive coefficient gap
$C-D_R$ is multiplied by the explicit positive function $1-w_*$.
A single uniform positive gap on the punctured leaf is unnecessary.
The weight records the precise way in which the upper bound can meet
the corner value. It thereby resolves the equality set as well as
proving a maximum bound.
On an ordinary leaf the largest validated coefficient sum is already
strictly less than $L_c$, giving the comparison everywhere on the leaf.

The partition is specified by a rooted binary tree: every internal node
splits its closed rectangle at the midpoint of one coordinate, and both
children must be present. Starting with $[-1,1]^2$, this reconstructs a
cover of the entire square. Shared edges belong to both adjacent leaves;
the same inequalities hold there. The designated corner vertices are
exactly the two nonnode corners in Theorem~\ref{thm:main-location}.
Their node-preserving reflection gives their common value $C$.
Thus the leaf predicates prove both the maximum and its exact equality
set on the full square.

The coverage statement can be checked by induction on the finite tree.
A node represents a closed rectangle $[u,v]\times[w,z]$. Splitting
its first coordinate at $m=(u+v)/2$ replaces it by
$[u,m]\times[w,z]$ and $[m,v]\times[w,z]$. Their union is the
parent, and they meet on the common vertical segment. A split in the
second coordinate has the analogous property. At a terminal node the
union of descendant leaves is just that node. Working upward through
the tree shows that the union of all leaves is the root square.
Every step is an exact equality of closed sets.

The checker therefore needs the tree structure together with the local
inequalities. Each internal node has both children, and the paths to
leaves form a complete collection with no pending region. Omitting a
child would invalidate the inductive union statement. For a complete
tree the number and shape of the leaves may vary substantially across
degrees while the coverage proof stays identical. A search can choose
where more subdivision is useful; the checker reconstructs the
specified final tree and verifies each resulting leaf.

Now fix an arbitrary physical point $z$. Coverage provides at least
one leaf containing $z$. An ordinary leaf gives
$\lambda_n(z)<C$. A corner leaf gives the same inequality unless
$z$ is its designated vertex. If $z$ is a designated vertex, the
cardinal formula supplies $\lambda_n(z)=C$. Shared boundaries cause
no difficulty: every applicable leaf bound is valid at the common
point, and any one strict bound suffices. Since the designated
vertices are the two nonnode corners, this argument proves precisely
the equality classification in Theorem~\ref{thm:main-location}.

All coefficient enclosures are reconstructed by the checker. The identities
\[
 T_{r+1}(x)=2xT_r(x)-T_{r-1}(x),\qquad
 s^r=\sum_{i=r}^n\frac{\binom{i}{r}}{\binom nr}B_i^n(s)
\]
reduce construction and basis conversion to exact integer and rational
operations once the cosine arguments are bracketed.
The power-to-Bernstein identity can also be checked without solving a
linear system. For $i\geq r$ the elementary binomial identity
$\binom ni\binom ir=\binom nr\binom{n-r}{i-r}$ gives
\[
 \begin{aligned}
 \sum_{i=r}^n\frac{\binom ir}{\binom nr}B_i^n(s)
 &=s^r\sum_{i=r}^n\binom{n-r}{i-r}
                      s^{i-r}(1-s)^{n-i}\\
 &=s^r\bigl(s+(1-s)\bigr)^{n-r}=s^r.
 \end{aligned}
\]
The change from a physical coordinate $x$ to $x=u+(v-u)s$ is an
ordinary binomial expansion with rational coefficients whenever the
leaf endpoints are rational. Each power $s^r$ is then converted by
this identity. Applying the conversion in both coordinates constructs
the tensor coefficients. All conversion weights are exact rational
numbers, so interval input coefficients can be propagated outwards.

Here is the midpoint subdivision behind the reference to de Casteljau.
For a one-variable degree-$n$ polynomial with Bernstein coefficients
$d_0,\ldots,d_n$, define a triangular array by
\[
 d_i^{(0)}=d_i,\qquad
 d_i^{(r)}=\frac{d_i^{(r-1)}+d_{i+1}^{(r-1)}}2,
 \quad 1\leq r\leq n,\quad 0\leq i\leq n-r.
\]
The left half of the original interval has Bernstein coefficients
$d_0^{(0)},d_0^{(1)},\ldots,d_0^{(n)}$; the right half has
$d_0^{(n)},d_1^{(n-1)},\ldots,d_n^{(0)}$. The two lists share the
value at the midpoint. This follows by repeatedly combining adjacent
terms in the Bernstein expansion using the binomial recurrence.
For instance the left coefficient in position $j$ is
\[
 d_0^{(j)}=2^{-j}\sum_{i=0}^j\binom ji d_i.
\]
Precisely, $p(s/2)=\sum_{j=0}^n d_0^{(j)}B_j^n(s)$ on
$0\leq s\leq1$; expanding the right-hand side verifies the identity.
The right list
follows by the same calculation after reversing the coordinate.
Thus the averaging array is an exact change of representation on the
two child intervals.

For degree two, the left coefficients are
$d_0,(d_0+d_1)/2,(d_0+2d_1+d_2)/4$, and the right coefficients are
$(d_0+2d_1+d_2)/4,(d_1+d_2)/2,d_2$. Every operation is an average.
Tensor subdivision applies these arrays to coefficient rows or columns,
according to the coordinate being split. A two-dimensional midpoint
tree is therefore handled by repeated one-dimensional exact identities.
The cardinal polynomial remains the same on the union of the children.

If exact coefficients are represented by intervals on an integer scale
$D$, an average of $[a/D,b/D]$ and $[c/D,e/D]$ is enclosed by
\[
 \left[\frac{\lfloor(a+c)/2\rfloor}{D},
       \frac{\lceil(b+e)/2\rceil}{D}\right].
\]
The lower endpoint rounds down and the upper endpoint rounds up. The
integer scale may stay fixed, while the interval continues to contain
the exact rational average. Induction over the averaging array and
then over tree depth proves enclosure of every child coefficient.
The identity is exact at the level of real polynomials; the stored
interval representation preserves it by containment.
 Alternating arctangent
series in Machin's identity bracket $\pi$; alternating cosine series and
the Chebyshev recurrence bracket the grid values. Integer operations round
outwards at every interval step. The checker reconstructs the whole
partition and its coefficients. Successful certificates are finite witnesses.
The accompanying execution record identifies which degrees were newly rerun
through the common implementation.

The transcendental constants enter only through enclosures of the
specified exact coefficients. For example, the alternating expansion
of $\arctan x$ at $x=1/5$ or $x=1/239$ brackets the value between
successive appropriate partial sums, with the first omitted term
controlling the error. Machin's identity combines these rational
brackets into a rational bracket for $\pi$. Alternating cosine
expansions then bracket the required basic cosines, and the Chebyshev
recurrence propagates the grid values. The exact endpoint values
$1,-1$, and applicable zero values can be retained as identities.
The corrected cardinal formula expresses the polynomial coefficients
using finite sums and products of these quantities.

For an interval $[a,b]$, an absolute-value upper enclosure is
$\max(|a|,|b|)$. A lower enclosure is zero when the interval contains
zero, and otherwise $\min(|a|,|b|)$. These elementary rules justify the
coefficient absolute sums and a lower enclosure of the common corner
value, even when some cardinal signs are unresolved. Interval products
use the minimum and maximum of the four endpoint products, and each
subsequent rounding preserves containment. Because every cardinal and
every basis transform uses finitely many operations, induction over
those operations proves the required coefficient enclosures.

This gives two finite inductions behind the small-degree proof. The
arithmetic induction establishes that the stored intervals contain the
exact polynomial coefficients on each leaf. The geometric induction
establishes that the leaves cover the square. Positive-basis inequalities
then turn these finite facts into the continuous conclusion. The
execution inventory records the actual successful degree instances of
this common scheme. In particular, the fresh independent checks at
$2\leq n\leq16$ and the reused accepted instances at
$17\leq n\leq40$ have their stated individual provenance in that
inventory.

The degree-one case is explicit. Its nodes
are $(1,1),(1,-1),(-1,0)$, with cardinals
$(1+x+2y)/4$, $(1+x-2y)/4$, and $(1-x)/2$. Hence
\[
 \lambda_1(x,y)=\max\{1,|y|+(1-x)/2\},
\]
whose maximum is two, attained only at $(-1,\pm1)$.

Indeed the first two cardinals have sum $(1+x)/2$ and difference $y$.
For real numbers $a,b$, $|a+b|+|a-b|=2\max(|a|,|b|)$, so their
absolute sum is $\max((1+x)/2,|y|)$. The third cardinal is
nonnegative on the square. Adding it gives the displayed expression.
Here $|y|\leq1$ and $(1-x)/2\leq1$. Their sum can reach two
exactly when both inequalities are equalities, namely $x=-1$ and
$|y|=1$. This verifies the separate degree-one case by a direct affine
computation.

\subsection{The phase certificate}
For $41\leq n\leq319$ the common envelope is enclosed on
$\mathcal T=\{(a,b):a,b\geq0,\ a+b\leq1\}$, with phases
$(\alpha,\beta)=(\pi a,\pi b)$ and $k=n/(n+1)$.

The variables $a,b$ are dimensionless phase coordinates. Thus an
increment $\Delta a$ means an increment $\pi\Delta a$ in the radian
phase $\alpha$, and the analogous convention holds for $b,\beta$.
They are distinct from the physical coordinates $x,y$ of the square.
The closed triangle $\mathcal T$ is a fundamental region for the
phase envelope under its folding symmetries. The earlier closed odd
charts connect every physical point to one of these phase points and
a physical chart center. At a noncorner center the reflection defect
will improve the envelope comparison. At an actual nonnode corner a
separate descent witness will supply strictness near its origin.

The envelope $S_n$ is a finite absolute sum of primitive kernels.
Its definition retains every corrected kernel term and endpoint
weight. In particular a coincidence of phase arguments is evaluated
through the finite kernel itself. This makes the envelope continuous
on the entire closed triangle, including axes and chart seams. The
rectangle interpolation bound is applied to its smooth signed family,
which is now written explicitly.
In the notation of \eqref{ul:primitive}, put
\[
 f_\varepsilon(a,b)=\frac1{2\mathsf N}
       \sum_{\Gamma^-}\varepsilon_{rs}H_{rs}(\pi a,\pi b).
\]
Then $S_n(\pi a,\pi b)=\max_\varepsilon f_\varepsilon(a,b)$,
and $|f_\varepsilon|\leq S_n$. Every signed board sum has frequencies
at most $\pi,\pi k$ in these coordinates. A global norm $M$
and upper bounds at the vertices therefore give
\begin{align}\label{cert:phase-bound}
 S_n(\pi a,\pi b)&\leq q_R(a,b),\\
 q_R&=I_RW+\frac{M\pi^2}{2}
 \{(a-a_0)(a_1-a)+k^2(b-b_0)(b_1-b)\},\notag\\
 \max_Rq_R&\leq V_R+\theta_RM,
 \quad V_R=\max_zW_z,\quad
 \theta_R=\frac{\pi^2}{8}\{(a_1-a_0)^2+k^2(b_1-b_0)^2\}.\notag
\end{align}

Here a frequency bound describes how fast an oscillating term can
vary. The largest absolute frequency in one coordinate is called the
bandwidth in that coordinate. A term $e^{i\xi a}$ has first derivative $i\xi e^{i\xi a}$
and second derivative $-\xi^2e^{i\xi a}$. A uniform frequency bound
can therefore turn a global uniform norm into a derivative norm.
For the present finite trigonometric polynomials this is the
trigonometric Bernstein inequality. The word Bernstein here refers
to the derivative inequality; the positive polynomial basis used in
\eqref{cert:bernstein} is a separate object with the same historical
name. Both ingredients have elementary finite formulations.

One can see the frequency scaling directly from a typical factor in
the kernel. The term
$\cos(r(j\pi+\alpha)/n)$ has $\alpha$ frequencies $\pm r/n$.
Since $r\leq n$, these have magnitude at most one. Replacing
$\alpha$ by $\pi a$ multiplies the bound by $\pi$. In the other
coordinate a factor with index $s\leq n$ has $\beta$ frequencies
$\pm s/(n+1)$, of magnitude at most $k$. Replacing $\beta$ by
$\pi b$ gives the bound $\pi k$. Finite sums, shifts by board
indices, and fixed signs all preserve these frequency limits.

For completeness, the norm form of the derivative inequality can be
obtained from a finite differentiation identity. Let $T$ be a
trigonometric polynomial of integer degree at most $d\geq1$, meaning
a finite sum $T(x)=\sum_{p=-d}^d a_p e^{ipx}$. Write
$\|T\|_\infty=\sup_{x\in\mathbb R}|T(x)|$ for its global uniform
norm, and set
\[
 \tau_j=\frac{(2j-1)\pi}{2d},\qquad
 c_j=\frac{(-1)^{j-1}}{4d\sin^2(\tau_j/2)},
 \qquad 1\leq j\leq2d.
\]
The identity and its coefficient sum are
\[
 T'(x)=\sum_{j=1}^{2d}c_jT(x+\tau_j),\qquad
 \sum_{j=1}^{2d}|c_j|=d.
\]
These statements concern finite sums. To verify the first, test each
basis function $e^{ipx}$, $-d\leq p\leq d$. Finite Fourier
orthogonality on the $2d$ half-shifted roots gives the multiplier
$ip$ for $-d+1\leq p\leq d$ from
\[
 c_j=\frac1{2d}\sum_{\ell=-d+1}^d
                  i\ell e^{-i\ell\tau_j}.
\]
The orthogonality here is the elementary finite geometric-sum identity
\[
 \frac1{2d}\sum_{j=1}^{2d}e^{i(p-\ell)\tau_j}
   =\begin{cases}1,&p=\ell,\\0,&p\ne\ell,\end{cases}
 \qquad -d+1\leq p,\ell\leq d.
\]
For unequal indices their difference has magnitude less than $2d$,
and the geometric sum has ratio different from one but numerator
zero. For equal indices every summand is one. Substituting the
formula for $c_j$ in $\sum_j c_j e^{ip\tau_j}$ therefore gives
$ip$ throughout the stated range.
Differentiating the finite geometric sum before substituting
$e^{-2di\tau_j}=-1$ gives, for $z=e^{-i\tau_j}$,
\[
 \sum_{\ell=-d+1}^d\ell z^\ell
           =\frac{2z^{1-d}}{(1-z)^2},\qquad z^{2d}=-1.
\]
Substitution in the preceding formula for $c_j$, together with
$e^{id\tau_j}=i(-1)^{j-1}$, evaluates it as the displayed real
coefficient $c_j$. The remaining endpoint frequency $p=-d$ follows
from $e^{-id\tau_j}=-e^{id\tau_j}$, so it has multiplier $-id$.
This checks every basis frequency and hence every $T$ by linearity.

For the absolute coefficient sum, the finite product for the equally
spaced sine zeros yields, after logarithmic differentiation twice,
\[
 \sum_{j=0}^{p-1}\csc^2\left(x+\frac{j\pi}{p}\right)
       =p^2\csc^2(px).
\]
Indeed the product of these $p$ sine factors is a nonzero constant
multiple of $\sin(px)$; the two logarithmic derivatives remove the
constant. With $p=2d$ and $x=\pi/(4d)$, the sum is $4d^2$.
Dividing by $4d$ gives $\sum_j|c_j|=d$ as asserted. Taking
absolute values in the differentiation identity proves
$\|T'\|_\infty\leq d\|T\|_\infty$. Apply it again to $T'$
to obtain $\|T''\|_\infty\leq d^2\|T\|_\infty$.
If the argument is rescaled so the highest frequency is $\omega$,
the same chain rule gives $\omega^2$ in the second derivative bound.
All norms in this calculation extend over the full real argument.

Apply this one-variable calculation with the other coordinate held
fixed. Each complete signed board sum obeys
$|f_\varepsilon|\leq S_n\leq M$ on all real phases once a global
norm $M$ is supplied. Its frequency limits therefore give the two
second-derivative bounds used in Lemma~\ref{cert:positive}. The
constant $M$ bounds a whole signed sum at once. Since the conclusion
is uniform in the sign vector, taking the maximum gives the envelope
bound in \eqref{cert:phase-bound}.

Trigonometric Bernstein inequalities supply
$|(f_\varepsilon)_{aa}|\leq\pi^2M$ and
$|(f_\varepsilon)_{bb}|\leq\pi^2k^2M$ from the global norm.
The final line of \eqref{cert:phase-bound} uses positivity of the
bilinear weights and the elementary bound
$(a-a_0)(a_1-a)\leq(a_1-a_0)^2/4$, and similarly for $b$.
Thus $V_R+\theta_RM$ is a simple certified upper for the whole leaf.

The function $q_R$ in that display is a concrete quadratic polynomial
in $a,b$. It keeps the four interpolated vertex upper bounds and adds
two nonnegative curvature allowances. Each allowance vanishes on the
two corresponding sides of the rectangle and is largest halfway
between them. Taking $V_R$ and the maximum of each coordinate product
produces the simpler constant upper bound on the final line. This
constant can be conservative, but it applies throughout the full
rectangle and consequently throughout any closed portion of the
rectangle lying in $\mathcal T$.

The geometry of a leaf and the validity of interpolation are distinct
points. A rectangle can cross the hypotenuse $a+b=1$. Its four
vertices are still legitimate phase evaluation points, because the
signed functions and their global norm are defined on the whole
real phase plane. Interpolating on the entire rectangle is therefore
valid before restricting the result to its triangle intersection.
The established symmetries can fold an outside vertex to a canonical
phase point for evaluation. The widths in $\theta_R$ remain the
widths of the original rectangle.
The complete midpoint tree covers the triangle; a rectangle
is discarded only if its lower-left coordinate sum exceeds one.
Tangencies and shared boundaries remain in the cover.

For a rectangle $R=[a_0,a_1]\times[b_0,b_1]$ inside $[0,1]^2$,
the minimum of $a+b$ is $a_0+b_0$. Its triangle intersection is
empty precisely when $a_0+b_0>1$. If equality holds, the lower-left
vertex lies on the hypotenuse and must be retained. The root square
contains $\mathcal T$, and every midpoint split preserves the closed
union of its two children. Intersecting these union identities with
$\mathcal T$ proves that all retained leaves cover the entire closed
triangle. This is the same finite geometric induction as for the
physical-square certificate.

The word \emph{global} in the norm premise is essential to the
preceding derivative argument. A bound at the four vertices of one
leaf controls the interpolation data on that leaf. The derivative
estimate also involves translated arguments in its finite
differentiation identity, so it uses a bound valid at all real
phases. The next absorption argument obtains such a norm from finite
data. Its starting point is the existence of the true continuous
supremum, a fact that precedes any numerical estimate.

The norm premise is itself certified. Let $M_*$ be the global real-phase
supremum of $S_n$. The folding symmetries established in
Section~\ref{sec:argmax} reduce this supremum to $\mathcal T$.
Compactness supplies a maximizing point, and every signed sum has
absolute value at most $M_*$.
The finiteness follows from elementary continuity and compactness:
$S_n$ is a finite sum of absolute values of trigonometric polynomials,
so it is continuous, and its folding symmetries move every real phase
pair into the compact triangle. Its supremum is therefore the maximum
of a continuous function on that triangle. This establishes a finite
number $M_*$, even while its value is unknown. The inequality
$|f_\varepsilon|\leq S_n\leq M_*$ follows directly from the signed
sum definition and holds before any leaf test is performed.
Consequently the derivative argument is
valid initially with this true, possibly unknown, norm. If all leaves
are numerical and $\theta_R<1$, a leaf containing that point gives
$M_*\leq V_R+\theta_RM_*$. Absorbing its last term and then taking
the largest possible leaf ratio yields
\begin{equation}\label{cert:tree-norm}
 M_*\leq\max_R\frac{V_R}{1-\theta_R}.
\end{equation}

The logic can be unpacked at a single maximizing point $z_*$. Choose
a leaf $R_*$ containing it. The interpolation estimate, applied with
the mathematically existing norm $M_*$, gives
\[
 M_*=S_n(z_*)\leq V_{R_*}+\theta_{R_*}M_*.
\]
When $\theta_{R_*}<1$, subtracting the last term and dividing by the
positive number $1-\theta_{R_*}$ gives
$M_*\leq V_{R_*}/(1-\theta_{R_*})$. Since the location of $z_*$
is unknown, the certificate takes the largest of these ratios over
all triangle-meeting leaves. All the ratios use finite vertex data
and known widths. Their maximum is consequently a numerical upper
bound for the true global norm.

For an illustrative absorption inequality, suppose a leaf containing
a maximizer gave $V_R=7$ and $\theta_R=1/8$. Then
$M_*\leq7+M_*/8$ implies $M_*\leq8$. The same algebra works for
arbitrary validated vertex values with $\theta_R<1$. This example
only illustrates the rearrangement; the retained certificates have
their own recorded inputs. A numerical upper $M$ can be rounded
upward from the largest leaf ratio and then used as an already
established constant in subsequent local tests.

The maximum over leaves also explains why all triangle-meeting leaves
must provide their numerical vertex premises in this construction.
Any one of them might contain $z_*$. A missing local premise would
leave open the corresponding possibility for the maximizing point.
An analytic comparison on a far region has a different form and a
different dependency, so the retained hybrid construction uses the
independent grid norm when such a region is represented analytically.
Alternatively the quarter-grid points
$E=\{(i/4,j/4):i,j\geq0,\ i+j\leq4\}$ and their upper maximum $V_E$
give
\begin{equation}\label{cert:grid-norm}
 M_*\leq\frac{V_E}{1-\pi^2/64}<\frac{784}{663}V_E.
\end{equation}
Indeed the ten quarter-squares cover $\mathcal T$; outside vertices fold
back to $E$ by $(a,b)\mapsto(1-a,1-b)$, a phase symmetry.
Each square has side $1/4$, so its Peano error is at most
$\pi^2(1+k^2)M_*/128\leq\pi^2M_*/64$. Hence
$M_*\leq V_E+\pi^2M_*/64$, proving \eqref{cert:grid-norm}; the
rational factor follows from $\pi<22/7$.
The quarter-grid construction can be visualized as four rows of
squares: four in the bottom row, then three, then two, then one.
In formulas the ten rectangles are
\[
 \left[\frac i4,\frac{i+1}4\right]
 \times\left[\frac j4,\frac{j+1}4\right],
 \qquad i,j\geq0,\quad i+j\leq3.
\]
Their closed intersections with $\mathcal T$ cover it. For a point
$(a,b)\in\mathcal T$, choose the interval immediately below a grid
endpoint whenever a coordinate is an endpoint, with the first interval
used for zero. The resulting indices are between zero and three and
have sum at most three. Thus the point belongs to one of the listed
rectangles, including on the hypotenuse and at grid vertices.

Every vertex inside the triangle belongs to $E$. An outside vertex
of one of these rectangles has coordinate sum $5/4$. The map
$(a,b)\mapsto(1-a,1-b)$ sends it to a quarter-grid point of sum
$3/4$, hence to $E$, and preserves the envelope by the established
phase symmetry. Thus the same finite upper maximum $V_E$ bounds all
four vertices of every quarter-square, including the outside ones.
There are fifteen points in the geometric set $E$; exact phase
symmetries can identify equivalent evaluation arguments in the
retained arithmetic representation.

With a side width $1/4$, the coordinate product bound is $1/64$.
Multiplying by $M_*\pi^2/2$ and adding the $k^2$ contribution
produces $M_*\pi^2(1+k^2)/128$. Since $0<k<1$, this is at most
$M_*\pi^2/64$. At a maximizing point the resulting inequality is
\[
 \left(1-\frac{\pi^2}{64}\right)M_*\leq V_E.
\]
The coefficient on the left is positive. Finally
$\pi^2/64<(22/7)^2/64=121/784$, so its reciprocal is smaller
than $784/663$. This proves the safe rational norm factor used in
\eqref{cert:grid-norm}. The finite grid provides the vertex data;
the interpolation allowance provides the continuum control between
those vertices.
Either construction gives a
validated $M\geq M_*$ to insert in every subsequent leaf test.
For \eqref{cert:tree-norm}, every leaf meeting the triangle must
supply its numerical vertex bounds. A tree using an analytic far leaf
uses the independent grid norm \eqref{cert:grid-norm}.

The two norm constructions thus have explicit sufficient hypotheses.
The tree construction uses a complete numerical cover with
$\theta_R<1$ on every triangle-meeting leaf. The grid construction
uses the fixed quarter-grid values and their proved folding
identities. Both start from the same existing true supremum, and
both yield a global upper suitable for the derivative inequalities.
The checker can select an allowed construction and reconstruct its
bound. Its later leaf tests then use that certified $M$ and the
validated corner lower endpoint $L_c$ as separate premises.

On a leaf the certified upper for $q_R$ must be below $L_c$ on its far
part. On its actual near part $\rho\leq1/2$, at a noncorner center,
the required inequality is $q_R-\delta_R<L_c$, with the appropriate
proved reflection floor $\delta_R$. The floor's radius is
$\min\{1/2,\pi(a_1+b_1)\}$; it is used only where its physical guard
holds. To make the guard explicit, write $p_-\leq\pi\leq p_+$ for the
positive rational bracket. Since $\rho=\pi(a+b)$ on $\mathcal T$,
closed rational proxies for the two parts are
\[
 R\cap\mathcal T\cap\{a+b\leq1/(2p_-)\},\qquad
 R\cap\mathcal T\cap\{a+b\geq1/(2p_+)\}.
\]
The first contains every actual near point, and the second every actual
far point.
Here ``proxy'' means a rationally described enclosing set on which a
sufficient comparison is checked. The exact physical division is
$\pi(a+b)\leq1/2$ versus $\pi(a+b)\geq1/2$. If an actual point
is near, then $p_-(a+b)\leq\pi(a+b)\leq1/2$, proving membership
in the first proxy. If it is far, then
$p_+(a+b)\geq\pi(a+b)\geq1/2$, proving membership in the second.
The containments follow because $a+b\geq0$ and the rational bracket
is positive. These two elementary multiplications determine the
direction of each proxy inequality.

Since $p_-\leq p_+$, one has
$1/(2p_+)\leq1/(2p_-)$. The two proxies therefore overlap in a
closed strip containing the actual seam. A point in this strip may
cause both sufficient comparisons to be required. The width of the
strip expresses the uncertainty in the bracket for $\pi$. Its
inclusion makes the arithmetic classification valid throughout the
closed domains, including every boundary point.
Their overlap retains the seam $\rho=1/2$. Requiring the
appropriate strict comparison on each nonempty proxy is sufficient;
the defect floor is invoked in the physical implication only at actual
near points. Its certified radius can be evaluated outwards as
$r_R=\min\{1/2,p_+(a_1+b_1)\}$ and substituted in
\eqref{ul:finite-floors}.
To verify this radius choice, take an actual near point in $R$.
Its physical radius obeys both $\rho\leq1/2$ and
$\rho=\pi(a+b)\leq p_+(a_1+b_1)$. It is therefore at most their
minimum $r_R$. The proved reflection-floor formula at that radius
bounds the defect for every eligible noncorner center and every
actual point with radius at most $r_R$. The floor has a physical
domain hypothesis, so its use in the final implication is tied to
the actual near case. The enclosing near proxy provides an
arithmetically convenient sufficient test for that case.

The separation between these sets is useful when reading the
certificate. The checker may establish an upper bound for $q_R$ on
the whole rectangle and use that stronger bound on its near and far
proxies. This is the simple constant $V_R+\theta_RM$ already proved
in \eqref{cert:phase-bound}. A rectangle is wholly near if an outward
upper on its phase radius is at most $1/2$; this too is a rational
inequality using $p_+$. Whichever admissible predicate is recorded,
the logical application at a physical point follows its actual
radius and its actual center.

The strict comparisons use rational enclosures throughout. For
example, $\pi^2$ in $\theta_R$ can be replaced by $p_+^2$,
increasing the curvature allowance. The factor $k=n/(n+1)$ and all
midpoint widths are rational. If $V_R,M,L_c$ are rational enclosures,
then the resulting upper bound and its difference from $L_c$ are
rational. The two proxy thresholds are rational as well. Thus deciding
which closed tests are required and checking their positive margins
can be done with integer numerators and positive denominators.
Every noncorner odd center has an eligible
projection, so \eqref{ul:defect} gives there
\[
 \lambda_n=S_n-\Delta_{JK}\leq q_R-\delta_n(r_R)<L_c
                                                   \leq\Lambda_c(n).
\]
On the actual far part $\lambda_n\leq S_n\leq q_R<L_c$.
Actual corners are covered by a direct descent witness.

The eligible projection used here is the one supplied by the earlier
reflection argument for a noncorner odd center. Its existence ensures
that the defect floor is available for every noncorner center in the
near case. The envelope bound, the projection floor and the identity
$\lambda_n=S_n-\Delta_{JK}$ then form the displayed chain of
inequalities. At far points the envelope upper bound by itself gives
the strict comparison. Both comparisons end below $L_c$, which is at
most the exact corner value. The direct corner descent addresses the
remaining near points whose chart center is an actual maximizing
corner.

The corner witness has a different geometric aim from a leaf with a
uniform strict margin. At an actual corner the function equals the
proposed maximum, and nearby values approach it. A useful witness
therefore proves that the loss is bounded below by a positive multiple
of $A^2+B^2$ in inward radian phases $A,B$. That function vanishes
exactly at the origin. Three actual function values, together with
the exact value at the origin and the derivative allowance, suffice
for this conclusion through squared-variable interpolation.

Here is the reason for the squared coordinates. For a complete signed
cardinal sum put
\[
 p_\varepsilon(A,B)=\sum_\ell\varepsilon_\ell
 L_\ell\bigl(-\cos(A/n),-\cos(B/(n+1))\bigr).
\]
Each such function is even separately in $A$ and $B$, because both
physical coordinates are cosine functions. In particular its first
and third derivatives in either inward phase vanish at zero. It is
also a finite trigonometric polynomial of bandwidth at most $1$ in
$A$ and $k$ in $B$. All real inward phases give physical coordinates
in the square, so
\[
 |p_\varepsilon(A,B)|\leq
 \lambda_n\bigl(-\cos(A/n),-\cos(B/(n+1))\bigr)\leq M
\]
holds globally. The last bound follows from the already certified
envelope norm and the chart inequality $\lambda_n\leq S_n$.
The trigonometric derivative estimate therefore applies through
fourth order to every fixed signed sum.

Set $u=A^2/2$, $v=B^2/2$, and write
$h_\varepsilon(u,v)=p_\varepsilon(\sqrt{2u},\sqrt{2v})$.
The even power series in each phase gives a smooth function of
$u,v$ up to the closed axes. For $A>0$, differentiation gives
\[
 (h_\varepsilon)_{uu}
       =\frac{A(p_\varepsilon)_{AA}-(p_\varepsilon)_A}{A^3}.
\]
To bound this expression at small $A$, integrate the numerator rather
than estimating its two terms separately. If
$N(A)=A(p_\varepsilon)_{AA}-(p_\varepsilon)_A$, then
$N'(A)=A(p_\varepsilon)_{AAA}$ and $N(0)=0$.
Evenness also gives $(p_\varepsilon)_{AAA}(0,B)=0$, so
\[
 \begin{aligned}
 (h_\varepsilon)_{uu}(u,v)
 &=\frac1{A^3}\int_0^A t\int_0^t
                    (p_\varepsilon)_{AAAA}(s,B)\,ds\,dt\\
 &=\int_0^1\!\int_0^1 t^2
          (p_\varepsilon)_{AAAA}(rtA,B)\,dr\,dt.
 \end{aligned}
\]
The last kernel is nonnegative and has total mass $1/3$. The fourth
derivative bound is $M$ in $A$ and $k^4M$ in $B$. Consequently
\[
 |(h_\varepsilon)_{uu}|\leq M/3,\qquad
 |(h_\varepsilon)_{vv}|\leq k^4M/3.
\]
The integral expression has a continuous value at $A=0$, namely
$(p_\varepsilon)_{AAAA}(0,B)/3$, so the first estimate includes
the axis. The second includes its axis by the same reasoning. This
explicit extension explains why squared interpolation is valid at
the origin as well as in the interior.

Apply Lemma~\ref{cert:positive} to these signed functions on a squared
rectangle $[0,U]\times[0,V]$. Its curvature allowance is
\[
 \frac M6\{u(U-u)+k^4v(V-v)\}.
\]
Suppose the vertex upper data have losses at least $cU,cV,c(U+V)$
from the exact corner value $C$ at the other three vertices. Their
bilinear interpolant is at most $C-c(u+v)$, because this affine
function has exactly those four comparison values. Since
$u(U-u)\leq Uu$ and $v(V-v)\leq Vv$ on the closed rectangle,
\[
 \frac M6\{u(U-u)+k^4v(V-v)\}
 \leq\frac M6\max\{U,k^4V\}(u+v).
\]
Combining these inequalities and using $u+v=(A^2+B^2)/2$ gives
\eqref{ul:squared-peano}. The common estimates hold for every
$p_\varepsilon$, and taking their finite maximum gives the actual
Lebesgue function. This derives the corner comparison from positive
interpolation while keeping all cardinal sign changes included.

The local descent uses \eqref{ul:squared-peano}. Take
$U=\pi^2/32$, $V=9\pi^2/512$, and outward actual-function values
$W_x,W_y,W_{xy}$ at $(\sqrt{2U},0),(0,\sqrt{2V}),
(\sqrt{2U},\sqrt{2V})$. With $u'=U/\pi^2$, $v'=V/\pi^2$, set
\[
 q=\min\left\{\frac{L_c-W_x}{u'},\frac{L_c-W_y}{v'},
                         \frac{L_c-W_{xy}}{u'+v'}\right\}.
\]
For the other three vertices the losses relative to the exact corner
value are at least $L_c-W_x$, $L_c-W_y$, and $L_c-W_{xy}$.
They are therefore at least $(q/\pi^2)U$, $(q/\pi^2)V$, and
$(q/\pi^2)(U+V)$, respectively. At the origin the value is exactly
$\Lambda_c(n)$. Apply \eqref{ul:squared-peano} with $c=q/\pi^2$
to obtain
\[
 \begin{gathered}
 \lambda_n(-\cos(A/n),-\cos(B/(n+1)))
 \leq\Lambda_c(n)-\frac{g}{2\pi^2}(A^2+B^2),\\
 g=q-\frac{M\pi^4}{6}\max\{u',k^4v'\}.
 \end{gathered}
\]
Thus the strict test $g>0$ proves positive quadratic loss off the
origin.
The normalization in this calculation can be checked one factor at a
time. The chosen dimensionless squared widths are
$u'=1/32$ and $v'=9/512$. The squared radian-phase widths are
$U=\pi^2u'$ and $V=\pi^2v'$, so the lower vertex loss
$q u'$ equals $(q/\pi^2)U$. Similarly $qv'=(q/\pi^2)V$, and
the diagonal loss is at least $(q/\pi^2)(U+V)$. Hence the affine
loss coefficient in squared radian coordinates is $c=q/\pi^2$.
The curvature cost in that coefficient is
$(M/6)\max(U,k^4V)=(M\pi^2/6)\max(u',k^4v')$.
Factoring out $1/\pi^2$ leaves precisely
\[
 c-\frac M6\max(U,k^4V)
   =\frac1{\pi^2}
       \left(q-\frac{M\pi^4}{6}\max(u',k^4v')\right)
   =\frac g{\pi^2}.
\]
The further factor $1/2$ comes from $u+v=(A^2+B^2)/2$.
This accounts for both powers of $\pi$ and the coefficient in the
displayed corner inequality.

Only upper actual values and a lower corner value are needed for
the three losses. For example, if $W_x$ bounds the actual value at
the $A$-axis test point, then
$C-\lambda_n(-\cos(\sqrt{2U}/n),-1)$ is at least
$C-W_x\geq L_c-W_x$. The analogous inequalities hold at the
other two points. Defining $q$ by the minimum makes all three
normalized losses hold at once. A positive accepted $g$ then implies
a positive $q$ and a strict decrease away from the origin. Every
arithmetic operation in this test can be made rational by replacing
$\pi^4$ in its cost by $p_+^4$. If the resulting rational lower
gap is positive, the true gap $g$ is positive as well.

The three actual test points use the physical inward angles
$\pi/(4n)$ and $3\pi/(16(n+1))$, with a zero angle in the other
coordinate for the axis tests. Equivalently these angles are the
rational multiples $1/(4n)$ and $3/(16(n+1))$ of $\pi$.
Indeed $\sqrt{2U}=\pi/4$ and $\sqrt{2V}=3\pi/16$ are radian
phases, and the physical cosine arguments divide those phases by
$n$ and $n+1$. Keeping the radian phases, normalized phases and
physical inward angles distinct explains the scaling used by the
point evaluations and by the interpolation rectangle.
The rectangle is $0\leq A\leq\pi/4$,
$0\leq B\leq3\pi/16$, which contains the inward half diamond
$A+B\leq1/2$. Complete signed cardinal sums are even in each inward
variable and have frequencies at most $1,k$; their absolute values are
bounded by the same global $M$, since $\lambda_n\leq S_n$ in every
chart. These are precisely the premises of the squared-variable
interpolation argument. The parity reflection transfers the witness to
the other actual corner.

The containment of the required near diamond is also explicit.
For nonnegative $A,B$ with $A+B\leq1/2$, both coordinates are at
most $1/2$. The endpoints $\pi/4$ and $3\pi/16$ both exceed
$1/2$; for the second it is enough that $\pi>8/3$.
Thus every inward near point belongs to the rectangle where the
quadratic loss holds. At $A=B=0$ the exact value is $C$.
For every other point of that rectangle the sum $A^2+B^2$ is
positive, giving the desired strictness. The parity-preserving
reflection permutes the cardinal polynomials and carries this
comparison to the other nonnode corner while keeping each coordinate
mesh in its original role.

From degree 256 the analytic estimate \eqref{ul:corner256} supplies
the same strict descent on the required diamond. An analytic far leaf
is permitted from degree 192 only when its entire triangle intersection
satisfies $\rho\geq\pi/2$. Equation~\eqref{ul:far192} then gives
$S_n<\Lambda_c(n)-3/100$ throughout that closed intersection.

These analytic substitutions preserve their full domain hypotheses.
The corner estimate from degree 256 applies on the inward square
specified in \eqref{ul:corner256}, and its positive quadratic loss
covers the required near diamond. The far estimate from degree 192
uses the radian guard $\rho\geq\pi/2$, equivalent on
$\mathcal T$ to $a+b\geq1/2$. This guard differs from the
near/far division at $\rho=1/2$ used in the numerical predicates.
The analytic far estimate can therefore replace a numerical far
leaf only when its entire relevant closed domain satisfies the
stronger half-$\pi$ guard. In the retained rectangle predicate, a
lower-left sum at least $1/2$ ensures that condition on the whole
rectangle and hence on its triangle intersection.

A rectangle that straddles this analytic guard continues to need the
allowed numerical comparison on the part left uncovered by an
analytic estimate. The closed guard includes its equality line.
Together with the closed tree coverage, this ensures that applying
the analytic far comparison creates no omitted boundary strip.
The available analytic estimates change the kind of witness stored
for eligible regions and degrees, while their proved implications
supply exactly the strict comparisons required by the global
argument.

To complete the implication, choose an odd chart for any point of the
physical square and fold its absolute phases into $\mathcal T$.
The tree gives a leaf containing those phases. Its far comparison,
noncorner near comparison, or actual-corner descent excludes the point
strictly, except when it is the origin of an actual corner chart.
Those origins are the two nonnode corners already identified by parity.
Their exact common value and the closed chart cover prove the asserted
maximum and equality set, including every chart tie and subdivision
boundary.

One can follow the conclusion at a fixed arbitrary point as follows.
First the closed odd-chart cover supplies a center and signed phases.
The folding identities supply a point of $\mathcal T$ with the
same envelope value. A complete tree supplies a leaf containing that
phase point. If an allowed full-leaf strict envelope bound applies,
then $\lambda_n\leq S_n<\Lambda_c(n)$ there for every center.
Otherwise the near and far predicates are applied according to the
actual radius of the chosen chart. The near noncorner case subtracts
the proved defect. The near actual-corner case uses the local
quadratic descent. At the seam either valid comparison retains its
closed boundary. These possibilities cover all points in the chart.

The only surviving equality possibility is zero inward phase in an
actual nonnode-corner chart. The two such physical corners have
value $\Lambda_c(n)$ by the exact cardinal identity and the
node-preserving reflection. A physical corner that is a Padua node
has Lebesgue value one by cardinal interpolation, and the closed
chart argument covers it with the other points. Chart ties merely
provide multiple valid ways to reach a comparison. All local bounds
include their closed edges, so every physical point receives the
same final strict exclusion or the identified exact corner value.

This completes the finite-to-continuous implication for the phase
certificates. The physical Bernstein certificates prove the identical
conclusion on $2\leq n\leq40$. The phase certificates handle
$41\leq n\leq319$, and the earlier analytic proof starts at
$n=320$. Thus the degree ranges meet consecutively, and the exact
equality classification is consistent across the finite and analytic
parts of Theorem~\ref{thm:main-location}.

\subsection{Evidence and arithmetic}
The retained proof inventory covers every integer $2\leq n\leq319$;
Section~\ref{sec:argmax} covers the rest. For the phase range all 280
primary/fresh pairs, including the control degree 511, were checked by
exact rational reconstruction of 560 essential proofs. Native outward
endpoints come from the frozen Arb/FLINT/MPFR/GMP implementation;
finite binary64 endpoints are converted losslessly to rationals for the
cover predicates. The arithmetic implementation, compiler and native
interfaces are trusted components of this computer-assisted proof.

There are several linked levels in this evidence statement. The
mathematical reduction specifies the meaning of a valid coefficient,
vertex enclosure, complete tree and strict leaf predicate. A concrete
witness specifies finite data for a particular degree. The checker
reconstructs its topology and rational consequences against the
specified formulas. The execution record connects those data to the
frozen arithmetic source and to the actual native evaluations. Each
level is needed to connect the theorem to its retained finite
instances.

The count for the phase range is transparent. There are
$319-41+1=279$ consecutive degrees from 41 through 319. The control
degree 511 supplies one additional degree, for 280 primary/fresh
pairs. Two retained essential proofs in each pair give 560 essential
proof reconstructions. The control is included in the checking
inventory; the all-degree coverage itself uses the consecutive finite
range and the analytic tail as described above. The small-degree
execution inventory has its separate fresh and reused provenance
within the common Bernstein scheme.

A finite binary64 number has the exact form $m2^e$ for integers
$m,e$. Once a native outward endpoint has been recorded in that
format, converting it to a rational preserves its exact represented
value. A negative exponent gives a power-of-two denominator. The
rational checker can then compute norms, proxy guards, fourth-power
corner costs and strict margins with exact arithmetic. This conversion
preserves the native enclosure endpoints. The correctness of the
outward endpoint computation itself is supplied by the frozen
arithmetic implementation and its trusted native dependencies.

For the reader, this explains which part of a comparison is a real
analysis statement and which part is an arithmetic check. Positive
interpolation proves that valid endpoints and derivative norms imply
a whole-region bound. Exact rational operations prove the strict
inequality between the recorded bounds. The native interval stage
supplies the endpoint containment assertions at its stated trust
boundary. The source identity and witness hashes connect those
assertions to the implementation and execution records being used.
The compiler and interfaces are included explicitly because they
participate in carrying that implementation to its recorded outputs.

The primary and fresh passes ran separately on the production machine.
Locally all essential proof/account/source links were checked. Complete
histories for two randomly selected degrees, 199 and 238, were inspected,
and fresh native verification was run at both degrees with distinct
review. These sample replays supplement
the checks of all essential proofs and links. Full unsampled
histories remain on the production machine under the selective-retention
policy. Source pins, witness identities, checker commands and the exact
new small-degree run inventory are in the companion reproducibility notes.
They permit regeneration; a durable public release remains to be arranged.
This evidence boundary is part of the theorem's computational scope.

The distinction between essential proofs and complete histories also
has a concrete purpose. An essential proof retains the finite inputs
and conclusions needed to reconstruct the certificate implication,
with its source and account links. A complete history records the
broader production execution around that witness. All essential
proofs in the stated inventory were checked. The detailed history
inspection and fresh native verification described here concern
the selected degrees 199 and 238, with their separate review. The
selective-retention policy specifies where the remaining complete
histories reside and how the retained links identify them.

These records make the computational premises inspectable at their
stated scope. The reproducibility notes identify the exact source
pins, witness identities and commands needed for regeneration, and
record which small-degree instances were newly checked. The public
release of a durable complete package remains an archival task.
The mathematical reading of the theorem keeps these evidence and
trust premises attached to the finite computations that support it.
The analytical corner expansion in the earlier theorem keeps its
own stated proof and dependencies.
